% 1- Explicar detenidamente la principal motivacion del articulo. Dejar claro cual es el posible problema
% 2- Presentar las preguntas de investigacion pertinentes
% 3- Presentar algunos resultados preliminares
% 4- Presentar el trabajo llevado hasta ahora de investigacion y desarrollo
% 5- Dejar algunas preguntas abiertas y problemas encontrados hasta ahora
%
% $Id: slides.tex 4228 2006-06-21 21:55:12Z jjamor $
%
%
% Compilar a .pdf con LaTeX (pdflatex)
% Es necesario instalar Beamer (paquete latex-beamer en Debian)
%

%
% Grficos:
% Los grficos pueden suministrarse en PNG, JPG, TIF, PDF, MPS
% Los EPS deben convertirse a PDF (usar epstopdf)
%


\documentclass{beamer}
\usetheme{Warsaw}
\usebackgroundtemplate{\includegraphics[width=\paperwidth]{format/libresoft-bg.png}}
\usepackage[english]{babel}
\usepackage[utf8]{inputenc}
\usepackage{graphics}
\usepackage{amssymb} % Simbolos matematicos
\usepackage{url}

%\definecolor{libresoftgreen}{RGB}{162,190,43}
%\definecolor{libresoftblue}{RGB}{0,98,143}

%\setbeamercolor{titlelike}{bg=libresoftgreen}

%% Metadatos del PDF.
\hypersetup{
  pdftitle={Libresoft},
  pdfauthor={Daniel Izquierdo Cort\'azar},
  pdfcreator={GSyC/Libresoft},
  pdfproducer=PDFLaTeX,
  pdfsubject={Exercise: metrics tools for FLOSS projects},
}
%%



\AtBeginSection[]
{
  \begin{frame}<presentation>
    \frametitle{Index}
    \tableofcontents[current]
  \end{frame}
}




\begin{document}



\title{Metrics tools for FLOSS projects
}
\subtitle{ Master on Free Software
}
\institute{jfelipe@libresoft.es\\
GSyC/Libresoft, Universidad Rey Juan Carlos}
\author{Felipe Ortega, Pedro Coca, Daniel Izquierdo Cort\'azar}
\date{Fuenlabrada, Madrid\\ 2nd, December 2010}

\frame{
\maketitle
\begin{center}
\includegraphics[width=4cm]{format/gsyc-urjc.pdf} 
\end{center}
}


% Si el titulo o el autor se quieren acortar para los pies de pgina
% se pueden redefinir aqu:
%\title{Titulo corto}
%\author{Autores abreviado}


%% LICENCIA DE REDISTRIBUCION DE LAS TRANSPAS
\frame{
~
\vspace{4cm}

\begin{flushright}
{\tiny
(cc) 2010 Felipe Ortega. \\
Some rights reserved. This document is distributed under the Creative \\
            Commons Attribution-ShareAlike 3.0 licence, available in \\
            http://creativecommons.org/licenses/by-sa/3.0/

%  Este documento (o uno muy similar) está disponible en \\
%  \url{http://gsyc.escet.urjc.es/~jjamor/}
}
\end{flushright}
}
%%

% \frame{
% \frametitle{Table of contents}
% \tableofcontents
% }


%\section{Introduction}


\begin{frame}
\frametitle{Motivation}
\begin{center}
\begin{itemize}
\item We need to gather metrics about FLOSS projects
\item Many different aspects to cover:
\begin{itemize}
 \item Source code metrics.
 \item Mailing list metrics.
 \item Bug tracking metrics.
 \item Licenses and documentation metrics.
 \item Development process metrics.
 \item (You name it)
\end{itemize}
\end{itemize}
\end{center}
\end{frame}

% \section{Activity}

\begin{frame}
\frametitle{Explore tools computing metrics from FLOSS projects}
\begin{center}
\begin{itemize}
 \item Explore the Internet to find metrics tools for FLOSS projects.
 \item Ideally, the complete set of tools should cover most aspects to be measured.
 \item Probably, tools will come from different projects/developers/communities.
 \item Some metrics will be provided by more than one tool. Which is the best one? Why?
 \item Which aspects are not covered at all by any available tool? Why?
\end{itemize}
\end{center}
\end{frame}

% \section{Results}

\begin{frame}
\frametitle{Results}
\begin{center}
\begin{itemize}
 \item Complete list of metrics tools, inclduing:
 \begin{itemize}
  \item Provider/autors.
  \item URL.
  \item Included metrics.
  \item Target information sources (SCM, BTS, mailing lists, forums...).
  \item Brief description.
  \item Case studies or examples.
 \end{itemize}

\end{itemize}
\end{center}
\end{frame}

\end{document}
